% GENERATED FILE. Edit canonical lessons, then run npm run generate:books.
% canonical-source-hash: fnv1a64:ff291f4c304f6922
% canonical-lessons: ZH-C526-cheng-weigh, ZH-C526-gongjin, ZH-C526-zhongliang, ZH-W526-liang-amount, ZH-C526-daliang, ZH-C526-shaoliang

\chapter{How Heavy?}
\label{ch:zh-how-heavy}
\hlchaptermodality{526}

\begin{chapteropening}
\textbf{By the end of this chapter:} \emph{I can weigh things, talk about kilograms and weight, and say I eat or drink a lot or a little.}

Complete the last lesson of chapter 526: I can weigh things, talk about kilograms and weight, and say I eat or drink a lot or a little.
\end{chapteropening}

\section[chēng]{chēng — to weigh on a scale}

\label{lesson:ZH-C526-cheng-weigh}



\begin{quote}
Before the new one: say the Mandarin for to make a formal complaint, then the Mandarin for to give the money back.
\end{quote}

\subsection*{You'll want to know: chēng}
\emph{chēng} — \textquotedblleft{}to weigh on a scale\textquotedblright{}.

\begin{quote}
\emph{qǐng nǐ bāng wǒ chēng zhèxiē píngguǒ} — Could you weigh these apples for me?
\end{quote}

Its characters come later; for now it is a word for the ear.

\begin{cousinweb}
\emph{chēng} is inherited from Middle Chinese \emph{tsyhing} (Guangyun, in Baxter's transcription), glossed 'to know light from heavy'; the dictionary also lists \emph{tsyhingH}, glossed 'fitting' and, in a common spelling, the balance that weighs things.
\end{cousinweb}

\begin{grammarlens}[title={What you've built}]
One more word for saying how it was.
\end{grammarlens}

\subsection*{Your turn}
\begin{itemize}
\raggedright
  \item \emph{Say it:} \emph{chēng}
  \item \emph{Say it:} \emph{chēng}, once more
  \item \emph{Say it:} \emph{qǐng nǐ bāng wǒ chēng zhèxiē píngguǒ} — Could you weigh these apples for me?
  \item \emph{Recall:} say the Mandarin for to make a formal complaint, then the Mandarin for to give the money back, then say \emph{chēng} again
\end{itemize}

\begin{tcolorbox}[breakable,colback=teal!4,colframe=teal!35!black,title={Before you move on}]
What does chēng mean? (\textquotedblleft{}To weigh on a scale\textquotedblright{}.) Say it once more.
\end{tcolorbox}
\section[gōngjīn]{gōngjīn — kilogram}

\label{lesson:ZH-C526-gongjin}



\begin{quote}
Before the new one: say the Mandarin for to give the money back, then the Mandarin for to weigh on a scale.
\end{quote}

\subsection*{You'll want to know: gōngjīn}
\emph{gōngjīn} — \textquotedblleft{}kilogram\textquotedblright{}.

\begin{quote}
\emph{wǒ bǐ qùnián zhòng le liǎng gōngjīn} — I'm two kilos heavier than last year.
\end{quote}

Its characters come later; for now it is a word for the ear.

\begin{cousinweb}
\emph{gōngjīn} is built from \zh{公} 'public, standard' and \emph{jīn}, a weight, both inherited from Middle Chinese, \emph{kuwng} and \emph{kj+n} (Guangyun, in Baxter's transcription), which glosses the second as sixteen ounces. The standard weight of a kilogram is a modern use.
\end{cousinweb}

\begin{grammarlens}[title={What you've built}]
One more word for saying how it was.
\end{grammarlens}

\subsection*{Your turn}
\begin{itemize}
\raggedright
  \item \emph{Say it:} \emph{gōngjīn}
  \item \emph{Say it:} \emph{gōngjīn}, once more
  \item \emph{Say it:} \emph{wǒ bǐ qùnián zhòng le liǎng gōngjīn} — I'm two kilos heavier than last year.
  \item \emph{Recall:} say the Mandarin for to give the money back, then the Mandarin for to weigh on a scale, then say \emph{gōngjīn} again
\end{itemize}

\begin{tcolorbox}[breakable,colback=teal!4,colframe=teal!35!black,title={Before you move on}]
What does gōngjīn mean? (\textquotedblleft{}Kilogram\textquotedblright{}.) Say it once more.
\end{tcolorbox}
\section[\texorpdfstring{\zh{重量} (zhòngliàng)}{重量 (zhòngliàng)}]{\zh{重量} (zhòngliàng) — weight; how heavy something is}

\label{lesson:ZH-C526-zhongliang}



\begin{quote}
Before the new one: say the Mandarin for to weigh on a scale, then the Mandarin for a kilogram.
\end{quote}

\subsection*{You'll want to know: \zh{重量}}
\textbf{\zh{重量}} — \emph{zhòngliàng} — \textquotedblleft{}weight; how heavy something is\textquotedblright{}.

\begin{quote}
\emph{zhège xiāngzi de zhòngliàng shì èrshí gōngjīn} — This suitcase weighs twenty kilos.
\end{quote}

The new character, in pieces: \zh{量} — twelve strokes in three pieces: \zh{日} on top, four strokes; a long bar under it; and \zh{里} at the foot, seven strokes. The next lesson puts it in your hand.

\begin{cousinweb}
\textbf{\zh{重量}} (\emph{zhòngliàng}) is built from \zh{重} 'heavy' and \emph{liàng} 'an amount', both inherited from Middle Chinese, \emph{drjowngX} and \emph{ljangH} (Guangyun, in Baxter's transcription), which glosses the first 'much, thick' and gives the second as units of dry measure. The dictionary lists the second again as \emph{ljang}, 'to measure', the liáng of the word for talking things over.
\end{cousinweb}

\begin{grammarlens}[title={What you've built}]
One more word for saying how it was.
\end{grammarlens}

\subsection*{Your turn}
\begin{itemize}
\raggedright
  \item \emph{Say it:} \emph{zhòngliàng}
  \item \emph{Say it:} \emph{zhòngliàng}, once more
  \item \emph{Say it:} \emph{zhège xiāngzi de zhòngliàng shì èrshí gōngjīn} — This suitcase weighs twenty kilos.
  \item \emph{Recall:} say the Mandarin for to weigh on a scale, then the Mandarin for a kilogram, then say \emph{zhòngliàng} again
\end{itemize}

\begin{tcolorbox}[breakable,colback=teal!4,colframe=teal!35!black,title={Before you move on}]
What does \zh{重量} mean? (\textquotedblleft{}Weight; how heavy something is\textquotedblright{}.) Say it once more.
\end{tcolorbox}
\section[\texorpdfstring{\zh{量} (liàng)}{量 (liàng)}]{\zh{量} — \zh{日}, a long bar, then \zh{里}}

\label{lesson:ZH-W526-liang-amount}



\begin{quote}
The word first, then the character inside it.

\begin{itemize}
\raggedright
  \item \emph{Recall:} say \emph{zhòngliàng} — \textbf{R1}, one lesson back, before it decays
\end{itemize}

\begin{quote}
\zh{里} \zh{日}
\end{quote}

\emph{Read it:} each one aloud, with nothing said first — \emph{lǐ}, \emph{rì}
\end{quote}

\subsection*{Script — \zh{日}, a long bar, then \zh{里}}
\hlblockfigure{\detokenize{figures/ZH-W526-liang-amount-filmstrip.pdf}}{How \zh{量} is written, stroke by stroke}

\begin{quote}
\zh{量}
\end{quote}

\textbf{Three pieces, top to bottom.} Count the pieces, not the strokes.

\textbf{The top piece, \zh{日}, four strokes}, small, as you wrote it: a vertical, a bar that turns down, a bar across the middle, a bar closing the foot.

\textbf{One long bar} across the middle, the widest stroke of the character.

\textbf{The bottom piece, \zh{里}, seven strokes}, as you wrote it: a vertical, a bar that turns down, a bar across the middle, a bar closing the box, a vertical down through the middle, a short bar, then the long bar at the foot.

4 + 1 + 7 = twelve.

\subsection*{Your turn}
\begin{itemize}
\raggedright
  \item \emph{Write it:} \zh{量}, all twelve, with the model in front of you
  \item \emph{Cover:} write it again from memory, and count the pieces, not the strokes
\end{itemize}

\begin{tcolorbox}[breakable,colback=teal!4,colframe=teal!35!black,title={Before you move on}]
How many pieces does \zh{量} have? (\textbf{Three: \zh{日}, a long bar and \zh{里}}.) How many strokes? (\textbf{Twelve}.)

Source: \href{https://raw.githubusercontent.com/chanind/hanzi-writer-data/68d10a4b21150cae5e1ebbd223eed289cf32d90c/data/\%E9\%87\%8F.json}{Hanzi Writer Data \zh{量}, pinned stroke paths and medians}.
\end{tcolorbox}
\section[\texorpdfstring{\zh{大量} (dàliàng)}{大量 (dàliàng)}]{\zh{大量} (dàliàng) — a large amount of; in large amounts}

\label{lesson:ZH-C526-daliang}



\begin{quote}
Before the new one: say the Mandarin for a kilogram, then the Mandarin for how heavy something is.
\end{quote}

\subsection*{You'll want to know: \zh{大量}}
\textbf{\zh{大量}} — \emph{dàliàng} — \textquotedblleft{}a large amount of; in large amounts\textquotedblright{}.

\begin{quote}
\textbf{\zh{今天格外热最好大量喝水}} \emph{jīntiān géwài rè zuìhǎo dàliàng hē shuǐ} — It's especially hot today, so you'd better drink plenty of water.
\end{quote}

\begin{cousinweb}
\textbf{\zh{大量}} (\emph{dàliàng}) is built from \zh{大} 'big' and \zh{量} 'an amount', both inherited from Middle Chinese, \emph{dajH} and \emph{ljangH} (Guangyun, in Baxter's transcription), which gives the second as units of dry measure.
\end{cousinweb}

\begin{grammarlens}[title={What you've built}]
One more word for saying how it was.
\end{grammarlens}

\subsection*{Your turn}
\begin{itemize}
\raggedright
  \item \emph{Say it:} \emph{dàliàng}
  \item \emph{Say it:} \emph{dàliàng}, once more
  \item \emph{Say it:} \emph{jīntiān géwài rè zuìhǎo dàliàng hē shuǐ} — It's especially hot today, so you'd better drink plenty of water.
  \item \emph{Recall:} say the Mandarin for a kilogram, then the Mandarin for how heavy something is, then say \emph{dàliàng} again
\end{itemize}

\begin{tcolorbox}[breakable,colback=teal!4,colframe=teal!35!black,title={Before you move on}]
What does \zh{大量} mean? (\textquotedblleft{}A large amount of; in large amounts\textquotedblright{}.) Say it once more.
\end{tcolorbox}
\section[\texorpdfstring{\zh{少量} (shǎoliàng)}{少量 (shǎoliàng)}]{\zh{少量} (shǎoliàng) — a small amount of}

\label{lesson:ZH-C526-shaoliang}



\begin{quote}
Before the new one: say the Mandarin for how heavy something is, then the Mandarin for in large amounts.
\end{quote}

\subsection*{You'll want to know: \zh{少量}}
\textbf{\zh{少量}} — \emph{shǎoliàng} — \textquotedblleft{}a small amount of\textquotedblright{}.

\begin{quote}
\textbf{\zh{吃了少量的米饭以后我感觉很好}} \emph{chī le shǎoliàng de mǐfàn yǐhòu wǒ gǎnjué hěn hǎo} — After a little rice I felt fine.
\end{quote}

\begin{cousinweb}
\textbf{\zh{少量}} (\emph{shǎoliàng}) is built from \zh{少} 'few' and \zh{量} 'an amount', both inherited from Middle Chinese, \emph{syewX} and \emph{ljangH} (Guangyun, in Baxter's transcription), which glosses the first 'not many'.
\end{cousinweb}

\begin{grammarlens}[title={What you've built}]
One more word for saying how it was.
\end{grammarlens}

\subsection*{Your turn}
\begin{itemize}
\raggedright
  \item \emph{Say it:} \emph{shǎoliàng}
  \item \emph{Say it:} \emph{shǎoliàng}, once more
  \item \emph{Say it:} \emph{chī le shǎoliàng de mǐfàn yǐhòu wǒ gǎnjué hěn hǎo} — After a little rice I felt fine.
  \item \emph{Recall:} say the Mandarin for how heavy something is, then the Mandarin for in large amounts, then say \emph{shǎoliàng} again
\end{itemize}

\begin{tcolorbox}[breakable,colback=teal!4,colframe=teal!35!black,title={Before you move on}]
What does \zh{少量} mean? (\textquotedblleft{}A small amount of\textquotedblright{}.) Say it once more.
\end{tcolorbox}
